\chapter{Présentation de l'entreprise}

\section{Histoire et données clés}

Winamax est un acteur français du numérique, spécialisé dans les plateformes de paris en ligne et de poker.
L'entreprise opère des services à forte disponibilité, avec des contraintes importantes de continuité de service,
de performance et de fiabilité.

Le stage s'inscrit dans ce contexte d'exploitation 24/7, au sein d'une organisation où la qualité du monitoring et
de l'alerting a un impact direct sur la capacité des équipes à détecter puis résoudre rapidement les incidents.

Les données chiffrées non confirmées (effectifs, volumes exacts, répartition d'équipes) ne sont pas détaillées dans
ce rapport à ce stade et seront complétées ultérieurement si nécessaire.

\section{Activité}

L'activité de Winamax repose sur des services applicatifs interconnectés, déployés dans une infrastructure cloud.
Ces services doivent rester disponibles en continu, notamment pendant les périodes de forte charge.

Dans ce cadre, l'astreinte technique joue un rôle central. Elle doit pouvoir distinguer rapidement :
\begin{itemize}
	\item les incidents réellement bloquants pour les utilisateurs,
	\item les signaux faibles ou anomalies ne nécessitant pas d'intervention immédiate.
\end{itemize}

Le sujet de stage répond précisément à cet enjeu, en visant une amélioration du système d'alerting pour réduire le
bruit opérationnel et mieux prioriser les interventions.

\section{Rayonnement international}

L'entreprise s'adresse à une base d'utilisateurs large et active, ce qui impose une exigence élevée de robustesse
des systèmes techniques et de réactivité des équipes d'exploitation.

Même sans détailler ici les marchés ou chiffres d'audience, ce contexte justifie l'importance d'un alerting orienté
impact utilisateur : une dégradation non détectée ou mal priorisée peut rapidement se traduire par des conséquences
fonctionnelles visibles à grande échelle.
